\documentclass[11pt]{article}
\usepackage{fullpage,url}
\usepackage{mathrsfs,amsmath}
\usepackage{graphicx}  
\usepackage{float}

\usepackage[letterpaper,top=1in,bottom=1in,left=1in,right=1in,nohead]{geometry}

\setlength{\parindent}{0in}
\setlength{\parskip}{6pt}

\DeclareMathOperator{\E}{E}
\DeclareMathOperator{\Var}{Var}
\DeclareMathOperator{\Unif}{Unif}

\begin{document}
\thispagestyle{empty}
{\large{\bf CSE 586A  \hfill Xiwen Li}}\\

{\LARGE{\bf Problem Set I}}
\vspace{0.2\baselineskip}
\hrule


\begin{enumerate}
%% Problem 1
\item The problem is to compare solution on a differential equation using two different differential equations.

Euler's approximation approximates both $v_{t}$ and $\phi_{t}$ in each iteration

\begin{enumerate}

	\item[(a)] 	$\frac{dy}{dt}+2y=2-e^{-4t}$\\
	\newline
	Let $\mu(t)=e(2t)$, $\mu'(t)=2e^{2t}$\\
	\newline
	Multiply both sides by $\mu(t)$: $\frac{de^{2t}y}{dt}+2e^{2t}y=2e^{2t}-e^{-2t}$\\
	\newline
	$\mu(t)\frac{dy}{dt}+\mu'(t)y=2e^{2t}-e^{-2t}$\\
	\newline
	$(\mu(t)y(t))'=2e^{2t}-e^{-2t}$\\
	\newline
	$y(t)=1+\frac{1}{2}e^{-4t}-\frac{C}{e^{2t}}$\\
	\newline
	Since $y(0)$=1,\\
	\newline
	$y(t)=1+\frac{1}{2}e^{-4t}-\frac{1}{2e^{2t}}$\\
	
	\item[(b)] We can rewrite the equation as: \\
	\newline
	$\frac{dy}{dt}=-2y+2-e^{-4t}$\\
	\newline
	Then the equation of Euler approximation is:\\
	\newline
	$d_{n+1}=d_{n}+h\frac{dy}{dt}$\\
	$d_{n+1}=d_{n}+h(-2y+2-e^{-4t})$
	\newline
	\begin{table}[h!]
		\begin{center}
			
			\label{tab:table1}
			\begin{tabular}{llllll} % <-- Alignments: 1st column left, 2nd middle and 3rd right, with vertical lines in between
				\textbf{step size} & \textbf{t = 1} & \textbf{t = 2}& \textbf{t = 3}& \textbf{t = 4}& \textbf{t = 5}\\
				\hline
				0.1 & 0.931324 & 0.991368 & 0.999050 & 0.999898 & 0.999989\\
				0.05 & 0.936470 & 0.991113 & 0.998898 & 0.999866 & 0.999984\\
				0.01 & 0.940499 & 0.991019 & 0.998789 & 0.999839 & 0.999979\\
				0.005 & 0.940996 & 0.991014 & 0.998776 & 0.999836 & 0.999978\\
				0.001 & 0.941391 & 0.991011 & 0.998766 & 0.999833 & 0.999977\\
				\hline
				original & 0.941490 & 0.991010 & 0.998764 & 0.999832 & 0.999977\\ 
				\hline
			\end{tabular}
			\caption{Values by Euler's Method}
		\end{center}
	\end{table}
	
	\item[(c)]
	
	
\end{enumerate}

%% Problem 2
\item 
The problem is to use Euler's approximation to estimate the transformation at t = 1.

At the very beginning, $v_{0}$ and $\phi_{0}$ are initialized. The initial image is defined by $I_{0}$. At each iteration, I approximate $v_{t}$ and $\phi_{t}$ using equation:

\begin{equation}
v_{t+1}=v_{t}+h*\frac{dv_{t}}{dt}
\end{equation}

\begin{equation}
\phi_{t+1}=\phi_{t}+h*\frac{d\phi_{t}}{dt}
\end{equation}

h is the step size. Once I get transformation, Final image is interpolated by as equation (3):

\begin{equation}
I_{final}=I_{0}\circ \phi_{t=1}
\end{equation}

\begin{enumerate}

\item
\begin{figure}[h]
	\begin{center}
		\includegraphics[width=0.6\textwidth]{./code/result_a}
		\caption{Final Image Transformed by (a)}
		\label{fig:f_res_a}
	\end{center}
\end{figure}


\item
\begin{figure}[h]
	\begin{center}
		\includegraphics[width=0.6\textwidth]{./code/result_b}
		\caption{Final Image Transformed by (b)}
		\label{fig:f_res_b}
	\end{center}
\end{figure}


\end{enumerate}


\end{enumerate}

\end{document}
